% 02-infrastructure.tex — Chapter 2: Infrastructure
\chapter{Infrastructure}
\label{chap:02-infrastructure}
\index{infrastructure}

\pnum
This chapter documents the utility layer that underpins all compiler subsystems.
The infrastructure provides primitive types, arena allocation, file I/O, ASCII
classification, string values, and the \cname{panic} assertion mechanism.  None of
these utilities depends on Lain-specific data; they are general-purpose C99 building
blocks found under \srcref{src/utils/}.

% -------------------------------------------------------------------------
\section{Primitive types}
\label{sec:infra-types}
\index{primitive types}
% -------------------------------------------------------------------------

\pnum
The file \srcref{src/utils/common/def.h} introduces fixed-width type aliases and
size-annotated arithmetic types used throughout the compiler.

\begin{center}
\begin{tabular}{lll}
\toprule
\textbf{Alias} & \textbf{Underlying C type} & \textbf{Usage in compiler} \\
\midrule
\cname{isize}  & \texttt{ptrdiff\_t} (signed)  & Array indices, offsets, counts \\
\cname{usize}  & \texttt{size\_t}   (unsigned) & Allocation sizes \\
\cname{iptr}   & \texttt{intptr\_t}            & Pointer arithmetic \\
\cname{uptr}   & \texttt{uintptr\_t}           & Alignment calculations \\
\cname{int8}\,\ldots\cname{int64}   & \texttt{int8\_t}\,\ldots\texttt{int64\_t}   & Fixed-width signed integers \\
\cname{uint8}\,\ldots\cname{uint64} & \texttt{uint8\_t}\,\ldots\texttt{uint64\_t} & Fixed-width unsigned integers \\
\cname{float32} & \texttt{float}  & (unused in current compiler) \\
\cname{float64} & \texttt{double} & (unused in current compiler) \\
\bottomrule
\end{tabular}
\end{center}

\pnum
The header also redefines \texttt{sizeof} and \texttt{alignof} to return
\cname{isize} and provides the \texttt{countof(a)} macro that divides the size of
an array by the size of its first element.

\begin{ccode}
#define sizeof(x)    (isize)sizeof(x)
#define alignof(x)   (isize)OPERATOR_ALIGNOF(x)
#define countof(a)   (sizeof(a) / sizeof(*(a)))
\end{ccode}

\begin{rationale}
Casting \texttt{sizeof} to \cname{isize} eliminates signed/unsigned comparison
warnings when loop indices (which are \cname{isize}) are compared to sizes.  This
is a compile-time cast; no runtime overhead is incurred.
\end{rationale}

% -------------------------------------------------------------------------
\section{Arena allocator}
\label{sec:infra-arena}
\index{arena allocator}
% -------------------------------------------------------------------------

\pnum
The arena allocator (\srcref{src/utils/arena.h}) is a minimal bump-pointer
allocator inspired by Chris Wellons's work.  It provides the memory foundation
for all three compiler arenas described in \cref{sec:arch-arenas}.

\pnum
The \cname{Arena} struct contains three \texttt{char *} pointers:

\begin{structdoc}{Arena}{src/utils/arena.h:23}
\cname{beg} & \texttt{char *} & Points to the first byte of the backing buffer. \\
\cname{cur} & \texttt{char *} & Points to the next free byte (the bump pointer). \\
\cname{end} & \texttt{char *} & Points one past the last byte of the buffer. \\
\end{structdoc}

\pnum
Allocation advances \cname{cur} by the requested number of bytes.  There is no
\texttt{free}; the only reclamation operation is \cname{arena\_clear()}, which
resets \cname{cur} to \cname{beg} and discards all previous allocations.

\begin{algorithm}
\textbf{arena\_push(arena, T)}
\begin{enumerate}
\item Assert \texttt{size > 0} and \texttt{count > 0}.
\item Assert \texttt{cur + size*count <= end} (sufficient space).
\item Save \texttt{ptr = cur}.
\item Advance \texttt{cur += size * count}.
\item Return \texttt{ptr} cast to \texttt{T *}.
\end{enumerate}
\end{algorithm}

\pnum
The aligned variant \cname{\_arena\_push\_aligned()} first computes the alignment
padding as:
\[
  \text{padding} = -((\text{uptr})\,\texttt{cur}) \mathbin{\&} (\text{alignment} - 1)
\]
This is the standard bitmask trick for aligning to a power-of-two boundary.  The
bump pointer is then advanced by \texttt{padding + size * count}.

\pnum
When \cname{ARENA\_DEBUG} is defined to a non-zero value (the default), every push
and alignment operation calls \cname{panic\_if\_msg()} to abort with a diagnostic if
any precondition is violated.  Debug mode has no effect on the generated code path
when all preconditions hold.

\begin{invariant}
\texttt{beg <= cur <= end} at all times.  No pointer returned by an arena push
is ever invalidated by a subsequent push, as the backing buffer is never
reallocated.
\end{invariant}

% -------------------------------------------------------------------------
\section{File I/O}
\label{sec:infra-file}
\index{file I/O}
% -------------------------------------------------------------------------

\pnum
Source files are read into \cname{file\_arena} by
\cname{file\_read\_into\_arena()} (\srcref{src/utils/file.h}).  This function
opens the file, queries its size, allocates \textit{size + 1} bytes from the
given arena, reads the entire file into that buffer, and NUL-terminates it.
The one extra byte allows all string-based consumers (the lexer in particular)
to treat the buffer as a C string without risk of reading past the end.

\pnum
The returned \cname{File} struct carries three fields:

\begin{structdoc}{File}{src/utils/file.h:17}
\cname{handle}   & \cname{file}   & Platform file handle (OS-specific typedef). \\
\cname{size}     & \cname{isize}  & Size of the file in bytes. \\
\cname{contents} & \cname{char *} & Pointer to arena-allocated NUL-terminated buffer. \\
\end{structdoc}

\pnum
If the file cannot be opened, \cname{contents} is set to \texttt{NULL} and
\cname{size} to zero.  The caller (\srcref{src/module.h}) checks for this
condition and emits a diagnostic before aborting.

% -------------------------------------------------------------------------
\section{Panic assertions}
\label{sec:infra-panic}
\index{panic assertions}
% -------------------------------------------------------------------------

\pnum
The file \srcref{src/utils/panic.h} provides a family of assertion macros that
print a structured error message to \texttt{stderr} and call \texttt{abort()}.
These are used exclusively for \textit{compiler-internal} invariants, not for
user-visible error reporting.

\begin{center}
\begin{tabular}{ll}
\toprule
\textbf{Macro} & \textbf{Behaviour} \\
\midrule
\cname{panic()}                   & Unconditional abort with file/line/function. \\
\cname{panic\_if(expr)}            & Abort if \texttt{expr} is true. \\
\cname{panic\_msg(fmt, ...)}       & Unconditional abort with formatted message. \\
\cname{panic\_if\_msg(expr,fmt,...)}& Abort if \texttt{expr} is true, with message. \\
\bottomrule
\end{tabular}
\end{center}

\pnum
The macro output format reports the source file, line number, and function name of
the failing assertion site.  This distinguishes internal compiler bugs from
user-visible semantic errors, which are reported via \cname{fprintf(stderr, \ldots)}
followed by \cname{exit(1)}.

% -------------------------------------------------------------------------
\section{ASCII classification}
\label{sec:infra-ascii}
\index{ASCII classification}
% -------------------------------------------------------------------------

\pnum
The header \srcref{src/utils/ascii.h} provides inline classification predicates
for ASCII characters.  The lexer relies on these predicates to classify input
bytes without depending on the locale-sensitive \texttt{<ctype.h>} functions,
which are undefined for non-ASCII byte values.

\pnum
The classifier functions include predicates for alphabetic characters
(\cname{ascii\_is\_alpha}), decimal digits (\cname{ascii\_is\_digit}),
hexadecimal digits, whitespace, and identifier-start / identifier-continue
characters.  Each is a \texttt{static inline} function that returns \texttt{bool}.

% -------------------------------------------------------------------------
\section{String type}
\label{sec:infra-string}
\index{string type}
% -------------------------------------------------------------------------

\pnum
The header \srcref{src/utils/string.h} defines a fat-pointer string struct:

\begin{ccode}
typedef struct {
    uint8 *data;
    isize  len;
} string;
\end{ccode}

\pnum
This type is used in utility-level code that operates on raw byte sequences
(hashing, comparison).  It is distinct from the compiler's own \cname{Id} type
(see \cref{sec:ast-identifiers}), which is an arena-allocated \cname{\{name, length\}}
pair used for identifier tokens.

\pnum
Two operations are provided: \cname{string\_hash(s)} computes a 64-bit hash of
the string content using a multiplicative hash, and \cname{string\_equals(a, b)}
compares two strings for equality using \cname{memcmp}.

\pnum
The \cname{string\_new(s)} macro wraps a C string literal:

\begin{ccode}
#define string_new(s)  { (uint8 *) s, lengthof(s) }
\end{ccode}

where \cname{lengthof(s)} is \cname{countof(s) - 1}, giving the string length
without the NUL terminator.
